\section{The positive-radius circle argument}
\label{sec:web-ckn-circle}

This section audits the incoming report's Section~13 (PDF pages 26--27,
printed pages 21--22). The geometric implication is valid. The proof below
specifies its solution class and proves the passage from a classical solution
before a terminal time to an interior singular set. It uses the
Caffarelli--Kohn--Nirenberg partial-regularity and suitable-existence theorems
as external theorems; it does not claim to reprove those theorems or to
certify the incoming report as a whole.

\subsection{The equation, suitable solutions, and the external theorem}

Let $\nu>0$ be the original viscosity. On an open set
$D\subset\mathbb R^3\times\mathbb R$, retain the equation
\begin{equation}\label{eq:ckn-original-ns}
 \partial_t u_i+\sum_{j=1}^3 u_j\partial_{x_j}u_i
 -\nu\sum_{j=1}^3\partial_{x_j}^2u_i+\partial_{x_i}p=f_i,
 \qquad \sum_{j=1}^3\partial_{x_j}u_j=0.
\end{equation}
Here $u:D\to\mathbb R^3$, $p:D\to\mathbb R$, and
$f:D\to\mathbb R^3$. Write $e=|u|^2/2$ and, for
$\phi\in C_c^\infty(D)$, define
\begin{equation}\label{eq:ckn-lei-functional}
 \mathcal L_\nu(u,p,f;\phi)=
 \int_D\left[e(\partial_t\phi+\nu\Delta\phi)
 +(e+p)u\cdot\nabla\phi+(f\cdot u)\phi
 -\nu|\nabla u|^2\phi\right]\,dx\,dt.
\end{equation}
The local energy inequality is $\mathcal L_\nu(u,p,f;\phi)\geq0$
for every nonnegative such $\phi$. In this section the sufficient suitable
class has, on each compactly contained cylinder $B\times I\Subset D$,
\begin{equation}\label{eq:ckn-integrability}
 u\in L^\infty(I;L^2(B))\cap L^2(I;H^1(B)),\qquad
 p\in L^{3/2}(B\times I),\qquad f\in L^q(B\times I),
 \quad q>5/2,
\end{equation}
and satisfies \eqref{eq:ckn-original-ns} distributionally and
\eqref{eq:ckn-lei-functional}. This is a sufficient class for the cited
CKN result; no assertion that these are minimal hypotheses is needed.
The products in the local energy inequality are integrable: the local
energy class embeds into $L^{10/3}$ in spacetime, hence into $L^3$ on a
bounded cylinder; $pu\in L^1$ follows from $p\in L^{3/2}$ and $u\in L^3$,
and $fu\in L^1$ follows from $q>5/2$ and the same finite-measure embeddings.
For clarity, the energy embedding follows from the spatial inequality
$\|u\|_{10/3}\leq C_B\|u\|_2^{2/5}\|u\|_{H^1}^{3/5}$:
raising to $10/3$ and integrating gives
$\int_I\|u\|_{10/3}^{10/3}\leq
C_B^{10/3}(\mathop{\rm ess\,sup}_I\|u\|_2^{4/3})
\int_I\|u\|_{H^1}^2$.

A regular point is a point with an open spacetime neighborhood on which
$u$ is essentially bounded. Denote the complementary set by $S(u)$.
Define
\begin{equation}\label{eq:ckn-parabolic-metric}
 d_{\rm par}((x,t),(y,s))=\max\{|x-y|,|t-s|^{1/2}\}.
\end{equation}
For a metric $d$, our Hausdorff measure convention is
\begin{equation}\label{eq:ckn-hmeasure}
 \mathcal H^1_d(E)=\lim_{\delta\downarrow0}
 \inf\left\{\sum_j\mathop{\rm diam}_d U_j:
 E\subset\bigcup_jU_j,\quad \mathop{\rm diam}_d U_j\leq\delta\right\}.
\end{equation}

\begin{theorem}[External CKN inputs]\label{thm:ckn-external}
For a suitable solution with divergence-free force locally in $L^q$,
$q>5/2$, the singular set has zero one-dimensional parabolic Hausdorff
measure. In the whole space, divergence-free $L^2$ initial data and a
divergence-free force in $L^2(\mathbb R^3\times(0,S))\cap L^q_{\rm loc}$,
$q>5/2$, admit a suitable energy weak solution on $(0,S)$ with the initial
trace, global energy inequality, and pressure in $L^{5/3}$ locally.
These are the inputs of CKN Theorem~B, Theorem~A', and Appendix~A;
the original paper uses unit viscosity. At regular points, its cited
interior regularity result gives bounded spatial derivatives on smaller
neighborhoods when the force has bounded smooth spatial derivatives.
\end{theorem}

The exact source and page checks are recorded in the accompanying source
inventory; the original publication is \cite{CKN82Verified}.
The circle consequence is also stated in the introduction of
Jiu--Xin \cite{JiuXin15Verified}, for their unforced equation. The forced
application here is obtained from the CKN hypotheses and the explicit
maps below, not by changing the force in Jiu--Xin's statement.
The readable research paper \cite{CLM2016Verified}, pp.2--3, independently
corroborates the force threshold, sufficient pressure class, and spatial
regularity assertion. The original CKN page images were unavailable;
that source-access limitation is retained in the inventory.
CKN's parabolic covering measure and \eqref{eq:ckn-hmeasure} have the same
null sets. For the direction used here it suffices that a parabolic
cylinder of radius $r$ has $d_{\rm par}$-diameter at most $2r$ (a temporal
interval of length $2r^2$ also has this bound). Every CKN cylinder cover
therefore yields a Hausdorff cover with diameter sum at most twice the
radius sum. Taking infima and then the small-radius limit proves that
CKN zero measure implies $\mathcal H^1_{d_{\rm par}}=0$.

\paragraph{Exact transfer of the external theorem for the original $\nu$.}
No viscosity is discarded. Introduce a second spacetime variable
$s=\nu t$ and the invertible map
\begin{equation}\label{eq:ckn-viscosity-map}
 \begin{aligned}
 D_\nu&=\{(x,\nu t):(x,t)\in D\},&
 v(x,s)&=\nu^{-1}u(x,s/\nu),\\
 \pi(x,s)&=\nu^{-2}p(x,s/\nu),&
 g(x,s)&=\nu^{-2}f(x,s/\nu).
 \end{aligned}
\end{equation}
The inverse is $u(x,t)=\nu v(x,\nu t)$,
$p(x,t)=\nu^2\pi(x,\nu t)$, $f(x,t)=\nu^2g(x,\nu t)$.
Each of $\partial_s v$, $(v\cdot\nabla)v$, $-\Delta v$,
$\nabla\pi$, and $g$ is respectively $\nu^{-2}$ times
$\partial_tu$, $(u\cdot\nabla)u$, $-\nu\Delta u$,
$\nabla p$, and $f$, evaluated at $(x,s/\nu)$. Also
$\nabla\cdot v=\nu^{-1}\nabla\cdot u$. Thus the equations,
including their distributional versions, are mapped exactly to the
unit-viscosity equations on $D_\nu$.
For $\psi(x,s)=\phi(x,s/\nu)$ and $ds=\nu\,dt$, direct substitution
in every term gives
\begin{equation}\label{eq:ckn-viscosity-lei}
 \mathcal L_1(v,\pi,g;\psi)
 =\nu^{-2}\mathcal L_\nu(u,p,f;\phi).
\end{equation}
For example, the time term contributes $\nu^{-2}$, the diffusion-test
term contributes $\nu^{-1}$ before comparison with its original factor
$\nu$, and the cubic, pressure, and force terms contribute
$\nu^{-3}\nu=\nu^{-2}$. The dissipation term contributes
$\nu^{-2}\nu=\nu^{-1}$ before comparison with its original factor
$\nu$. All coefficients therefore agree with
\eqref{eq:ckn-viscosity-lei}. The spacetime $L^q$ force integral is
multiplied by $\nu^{1-2q}$, the pressure $L^{3/2}$ integral by
$\nu^{-2}$, the essential supremum of the spatial velocity $L^2$
norm squared by $\nu^{-2}$, and the spacetime gradient $L^2$ integral
by $\nu^{-1}$. These positive finite factors preserve the stated classes.
The same substitution transfers the global energy inequality and initial
trace.
Finally, for $F_\nu(x,t)=(x,\nu t)$,
\begin{equation}\label{eq:ckn-metric-transfer}
 \min\{1,\sqrt\nu\}\,d_{\rm par}(a,b)
 \leq d_{\rm par}(F_\nu a,F_\nu b)
 \leq\max\{1,\sqrt\nu\}\,d_{\rm par}(a,b).
\end{equation}
Applying these inequalities to every covering set proves preservation of
Hausdorff null sets. Boundedness neighborhoods correspond under $F_\nu$
and multiplication of velocity by $\nu^{-1}$. Consequently the external
theorem gives its claimed zero-measure and existence conclusions for
the original equation with its original positive $\nu$.

\subsection{Rotation as an exact action on the equation and energy inequality}

For $\theta\in\mathbb R$, let
\begin{equation}\label{eq:ckn-rotation}
 \begin{gathered}
 R_\theta=\begin{pmatrix}
 \cos\theta&-\sin\theta&0\\
 \sin\theta&\cos\theta&0\\0&0&1
 \end{pmatrix},\\
 \mathcal R_Ru(x,t)=R u(R^Tx,t),\quad
 \mathcal R_Rp(x,t)=p(R^Tx,t),\quad
 \mathcal R_Rf(x,t)=R f(R^Tx,t).
 \end{gathered}
\end{equation}
These maps have domain the fields on $D$ and codomain the corresponding
fields on $RD=\{(Rx,t):(x,t)\in D\}$, and inverse $\mathcal R_{R^T}$.
For two rotations $R,Q$,
\[
 \mathcal R_R\mathcal R_Qu(x,t)=RQ\,u(Q^TR^Tx,t)
 =\mathcal R_{RQ}u(x,t).
\]
The scalar formula is identical without the prefactor. This proves the
group law, including the domains and inverses.

Write $y=R^Tx$. With every index ranging from 1 to 3, the derivative
identities are
\begin{align}\label{eq:ckn-rotation-components}
 \partial_t(\mathcal R_Ru)_i&=\sum_aR_{ia}\partial_tu_a(y,t),\nonumber\\
 \partial_{x_j}(\mathcal R_Ru)_i
 &=\sum_{a,b}R_{ia}R_{jb}\partial_{y_b}u_a(y,t),\nonumber\\
 \sum_i\partial_{x_i}(\mathcal R_Ru)_i
 &=\sum_a\partial_{y_a}u_a(y,t),\nonumber\\
 \Delta_x(\mathcal R_Ru)_i&=\sum_aR_{ia}\Delta_yu_a(y,t),\nonumber\\
 \sum_j(\mathcal R_Ru)_j\partial_{x_j}(\mathcal R_Ru)_i
 &=\sum_{a,b}R_{ia}u_b\partial_{y_b}u_a(y,t),\nonumber\\
 \partial_{x_i}(\mathcal R_Rp)&=\sum_bR_{ib}\partial_{y_b}p(y,t).
\end{align}
For divergence and Laplacian, use $\sum_iR_{ia}R_{ib}=\delta_{ab}$;
for the nonlinear term, insert
$(\mathcal R_Ru)_j=\sum_cR_{jc}u_c$ and contract
$\sum_jR_{jc}R_{jb}=\delta_{cb}$. Thus the full momentum residual
transforms to $R$ times the original residual at $(R^Tx,t)$, with the
same $\nu$, and the divergence constraint is preserved.
These calculations remain valid in distributions: in the defining
integrals for distributional derivatives make the linear substitution
$x=Ry$, whose Jacobian is $\det R=1$, and use the displayed chain rule
on each compactly supported test function. Tensor products transform
by $u^R\otimes u^R=R(u\otimes u)R^T$, so the distributional
divergence form of convection has the same identity.

Moreover
\begin{equation}\label{eq:ckn-rotation-norms}
 |u^R(x,t)|=|u(y,t)|,\qquad
 \nabla u^R(x,t)=R\nabla u(y,t)R^T,\qquad
 |\nabla u^R(x,t)|^2=|\nabla u(y,t)|^2.
\end{equation}
The last equality follows by expanding the Frobenius norm and contracting
both orthogonal factors. For $\phi\in C_c^\infty(RD)$ put
$\psi(y,t)=\phi(Ry,t)$. Then
$\nabla_y\psi=R^T\nabla_x\phi$, $\Delta_y\psi=\Delta_x\phi$,
and $\partial_t\psi=\partial_t\phi$ at corresponding points.
The force contraction is $(Rf)\cdot(Ru)=f\cdot u$ and the transport
contraction is $(Ru)\cdot\nabla_x\phi=u\cdot\nabla_y\psi$.
Consequently every integral in \eqref{eq:ckn-lei-functional} satisfies
\begin{equation}\label{eq:ckn-rotation-lei}
 \mathcal L_\nu(u^R,p^R,f^R;\phi)
 =\mathcal L_\nu(u,p,f;\psi).
\end{equation}
Nonnegative tests correspond bijectively, proving exact preservation of
suitability, the force exponent, and all local energy norms.
For completeness, the vorticity action has the same orientation:
\begin{equation}\label{eq:ckn-rotation-curl}
 \nabla_x\times u^R(x,t)=R(\nabla_y\times u)(y,t).
\end{equation}
Indeed the alternating trilinear form obeys
$\det(Ra,Rb,Rc)=\det(R)\det(a,b,c)=\det(a,b,c)$;
expanding this identity in coordinate vectors gives the alternating-tensor
contraction needed when \eqref{eq:ckn-rotation-components} is inserted
in $(\nabla\times u^R)_i=\sum_{j,k}\epsilon_{ijk}\partial_j u^R_k$.
The sign in \eqref{eq:ckn-rotation-curl} is therefore positive.

Axisymmetry means precisely
\begin{equation}\label{eq:ckn-axisymmetry}
 u(R_\theta x,t)=R_\theta u(x,t)
 \quad\hbox{for every $\theta$, as an equality of measurable fields.}
\end{equation}
It is vector equivariance, not invariance of each Cartesian component.
If the full triple is axisymmetric, it is a fixed point of the action.
For the singular-set argument only \eqref{eq:ckn-axisymmetry} is needed.
For each fixed $\theta$, the equality holds outside a null set; no common
null set for uncountably many angles is required, because essential
boundedness is compared separately for each fixed rotation.

\subsection{The fixed-time circle has positive parabolic measure}

Fix $r_0>0$, $z_0\in\mathbb R$, and a time $T$. Define the actual orbit
\begin{equation}\label{eq:ckn-circle}
 \mathcal C_{r_0,z_0,T}=
 \{(r_0\cos\theta,r_0\sin\theta,z_0,T):0\leq\theta<2\pi\}.
\end{equation}
The embedding $\iota_T:x\mapsto(x,T)$ is an isometry from Euclidean
space onto the time slice for \eqref{eq:ckn-parabolic-metric}.
This also gives equality of the corresponding Hausdorff measures for
subsets of that slice: intersecting an ambient cover with the slice
does not increase diameter, while lifting a spatial cover preserves
diameter. These operations prove the two inequalities between the
covering infima in \eqref{eq:ckn-hmeasure} for every $\delta$, hence
their equality after the limit.

Here is a direct positive lower bound that needs no curve-measure theorem.
The projection $(x,t)\mapsto x_1$ has Lipschitz constant one for
$d_{\rm par}$, and maps \eqref{eq:ckn-circle} onto $[-r_0,r_0]$.
If $(U_j)$ is any countable cover of that circle, the projections of
$U_j\cap\mathcal C_{r_0,z_0,T}$ can be enclosed in intervals of length
at most $\mathop{\rm diam}_{d_{\rm par}}U_j+\varepsilon2^{-j}$.
Those intervals cover $[-r_0,r_0]$. Countable subadditivity of interval
length therefore yields
$2r_0\leq\sum_j\mathop{\rm diam}_{d_{\rm par}}U_j+\varepsilon$.
Letting $\varepsilon\downarrow0$, taking the covering infimum and then
$\delta\downarrow0$ proves
\begin{equation}\label{eq:ckn-circle-positive}
 \mathcal H^1_{d_{\rm par}}(\mathcal C_{r_0,z_0,T})\geq2r_0>0.
\end{equation}
The measure is finite as well: partition $[0,2\pi]$ into intervals of
length at most $\delta/r_0$. The corresponding arcs cover the circle,
each arc has diameter no greater than $r_0$ times the interval length
by integration of the derivative of the circle parametrization, and
the total diameter sum is at most $2\pi r_0$. Thus the precise bounds
$2r_0\leq\mathcal H^1_{d_{\rm par}}(\mathcal C)\leq2\pi r_0$
are sufficient for the contradiction; no circumference normalization
of Hausdorff measure is silently imposed.

\begin{proposition}[Interior suitable circle exclusion]\label{prop:ckn-interior}
Let $D$ be rotationally invariant and let $(u,p,f)$ belong to the suitable
class to which Theorem~\ref{thm:ckn-external} applies. If $u$ is
axisymmetric, every point of $S(u)$ has $x_1=x_2=0$.
\end{proposition}
\begin{proof}
An orthogonal rotation maps any neighborhood of $(x,t)$ bijectively onto
a neighborhood of $(Rx,t)$ and preserves $|u|$ and spacetime null sets.
Therefore \eqref{eq:ckn-axisymmetry} says that either both points are
regular or both are singular. If a singular point has radius $r_0>0$,
every point of its orbit \eqref{eq:ckn-circle} is singular. Monotonicity
of the covering definition and \eqref{eq:ckn-circle-positive} imply
$\mathcal H^1_{d_{\rm par}}(S(u))\geq2r_0>0$, contrary to
Theorem~\ref{thm:ckn-external}. If the domain is not the whole space,
the complete orbit must be contained in the open domain, as ensured
here by rotational invariance; this is an interior argument.
\end{proof}

\subsection{The prescribed force need not be divergence-free}

For the whole-space application, take the original smooth force $f$
on $\mathbb R^3\times[0,S]$, with its original spatial decay. Define
\begin{equation}\label{eq:ckn-force-projection}
 \chi=\Delta^{-1}\nabla\cdot f
 =-\frac1{4\pi}\int_{\mathbb R^3}
 \frac{\nabla_y\cdot f(y,t)}{|x-y|}\,dy,
 \qquad g=f-\nabla\chi,\qquad \pi=p-\chi.
\end{equation}
The sign follows from the fundamental solution
$\Delta(-1/(4\pi|x|))=\delta_0$. To check that identity, the function
is harmonic away from zero and its outward radial flux through a sphere
of radius $\epsilon$ is
$4\pi\epsilon^2(1/(4\pi\epsilon^2))=1$; integration by parts against
a test function and passage to $\epsilon\downarrow0$ gives $\delta_0$.
Consequently $\Delta\chi=\nabla\cdot f$ and $\nabla\cdot g=0$.
Substituting \eqref{eq:ckn-force-projection} gives
\begin{equation}\label{eq:ckn-projected-equation}
 \partial_tu+(u\cdot\nabla)u-\nu\Delta u+\nabla\pi=g.
\end{equation}
This is an exact equivalence: the inverse restores $p=\pi+\chi$ and
$f=g+\nabla\chi$. It changes neither $u$ nor the original force being
prescribed in \eqref{eq:ckn-original-ns}.

The local energy functional is unchanged. The difference of the
projected and original integrands in the pressure and force terms is
$-\chi u\cdot\nabla\phi-(\nabla\chi\cdot u)\phi
=-u\cdot\nabla(\chi\phi)$, whose integral vanishes by
$\nabla\cdot u=0$. Hence
\begin{equation}\label{eq:ckn-force-projection-lei}
 \mathcal L_\nu(u,\pi,g;\phi)
 =\mathcal L_\nu(u,p,f;\phi).
\end{equation}
All distributional products here are legitimate because $\chi$ is smooth
on each compact spacetime set. To see this smoothness without estimating
a singular differentiated kernel, differentiate the smooth rapidly
decaying function $\nabla\cdot f$ under convolution with $|x-y|^{-1}$,
which is locally integrable; uniform decay bounds dominate the tail.

There is also a direct global bound needed for suitable existence.
For spatial Fourier frequency $\xi\ne0$, the multiplier of
$f\mapsto g$ is
$P(\xi)=I-\xi\xi^T/|\xi|^2$. Its square equals itself, its transpose
equals itself, and
$|P(\xi)a|^2=|a|^2-|\xi\cdot a|^2/|\xi|^2\leq|a|^2$.
Define its value at zero arbitrarily, since that set has zero measure.
Plancherel's identity proves, for each integer $m\geq0$,
\begin{equation}\label{eq:ckn-force-projection-norms}
 \|g(t)\|_{H^m(\mathbb R^3)}\leq\|f(t)\|_{H^m(\mathbb R^3)},
 \qquad
 \|g\|_{L^2(\mathbb R^3\times(0,S))}
 \leq\|f\|_{L^2(\mathbb R^3\times(0,S))}.
\end{equation}
Smooth force decay uniformly on this finite time interval makes every
right-hand side finite, and the local smoothness just proved gives
$g\in L^q_{\rm loc}$ for every finite $q$. Thus the original force,
even when not divergence-free, leads exactly to the force class of
Theorem~\ref{thm:ckn-external}. Adding $\chi$ to the resulting pressure
preserves local $L^{3/2}$ and restores the original equation and local
energy inequality. In the global energy inequality the work is also
unchanged: $g=Pf$, $Pv=v$ for divergence-free $L^2$ velocity $v$, and
self-adjointness of the matrix multiplier gives
$\langle g,v\rangle=\langle Pf,v\rangle=\langle f,Pv\rangle
=\langle f,v\rangle$.

\subsection{Existence beyond the terminal time and agreement before it}

Fix $0<T<S<\infty$, divergence-free smooth decaying data $u^0$, and
an original force $f$ smooth on $\mathbb R^3\times[0,S]$ with spatial
decay and derivative bounds there, for example the force class in the
official whole-space problem statement. Suppose $u$ is its classical
strong solution on $[0,T)$, in $C([0,t_0];H^m)$ for every $t_0<T$
and some integer $m\geq3$. This is the usual smooth strong-solution
branch from these data. These hypotheses specify the branch; mere
pointwise smoothness with $\|u(t)\|_2<\infty$ at each time is not being
silently substituted for them.

For every $t<T$, the classical global energy equality is
\begin{equation}\label{eq:ckn-classical-energy}
 \frac12\|u(t)\|_2^2+\nu\int_0^t\|\nabla u(s)\|_2^2\,ds
 =\frac12\|u^0\|_2^2+\int_0^t\langle f(s),u(s)\rangle\,ds.
\end{equation}
To derive it, pair \eqref{eq:ckn-original-ns} with $u$.
The convection term is $\int\nabla\cdot(|u|^2u/2)=0$,
the pressure term is $\int\nabla\cdot(pu)=0$ after the standard
whole-space cutoff limit, and integration by parts makes
$-\nu\int u\cdot\Delta u=\nu\int|\nabla u|^2$.
The strong Sobolev regularity on every shorter interval justifies these
limits and time integration; equivalently one may pair the projected
equation, whose pressure term vanishes by the orthogonal projection.
Set $F(t)=\int_0^t\|f(s)\|_2\,ds$ and $M(t)=\|u^0\|_2+F(t)$.
Regularizing the square root of $\|u\|_2^2$ by
$(\|u\|_2^2+\epsilon^2)^{1/2}$ in its differential inequality gives
derivative at most $\|f\|_2$. Integration and
$\epsilon\downarrow0$ yield
\begin{equation}\label{eq:ckn-uniform-energy}
 \|u(t)\|_2\leq M(t)\leq M(T),\qquad
 \nu\int_0^T\|\nabla u(s)\|_2^2\,ds\leq\frac12M(T)^2.
\end{equation}
For the second inequality insert $\|u(s)\|_2\leq M(s)$ into
\eqref{eq:ckn-classical-energy}, use
$\int_0^t F'(s)M(s)\,ds=(M(t)^2-M(0)^2)/2$, discard the nonnegative
terminal energy, and pass monotonically to $T$.
This proves the uniform energy and dissipation bounds from the specified
strong branch and force control; they are not inferred from finiteness
of individual time-slice energies alone.

Theorem~\ref{thm:ckn-external}, the force map
\eqref{eq:ckn-force-projection}, and the viscosity map
\eqref{eq:ckn-viscosity-map} give a suitable energy weak solution
$(v,P)$ for the original data and original force on $(0,S)$.
The suspected time $T$ is now an interior time. We next prove the needed
agreement; uniqueness after $T$ is neither used nor asserted.

Fix $t_0<T$. Put $w=v-u$ and let
$\kappa(t)=\|\nabla u(t)\|_{L^\infty}$, where the matrix norm is the
Frobenius norm. The assumed strong branch gives
$\int_0^{t_0}\kappa(t)\,dt<\infty$ by the $H^m$ embedding for $m\geq3$.
Testing the weak equation of $v$ with the smooth divergence-free $u$
and using the strong equation of $u$ gives the cross identity
\begin{align}\label{eq:ckn-cross-identity}
 \langle v(t),u(t)\rangle-\|u^0\|_2^2
 =\int_0^t\int_{\mathbb R^3}
 \bigg[\sum_{i,j}v_i(v_j-u_j)\partial_j u_i
 -2\nu\nabla v:\nabla u+f\cdot(u+v)\bigg]\,dx\,ds.
\end{align}
Here the term $\sum v_iv_j\partial_j u_i$ is the weak convection
term; the strong convection in $\partial_tu$ gives
$-\sum v_i u_j\partial_j u_i$; and the two diffusion pairings each
give $-\nu\nabla v:\nabla u$. This identifies every sign and factor.
The test is justified by time smoothing and divergence-free spatial
approximation of $u$: the convection pairing is bounded by
$\kappa\|v\|_2^2$, the diffusion pairing by
$\|\nabla v\|_2\|\nabla u\|_2$, and the force pairings by their
$L^2$ norms, so the approximations converge. The weak initial trace
and strong initial trace give the displayed initial cross term.

Add the energy inequality for $v$ to
\eqref{eq:ckn-classical-energy}, then subtract
\eqref{eq:ckn-cross-identity}. Expanding the norm squares yields,
for almost every $t\leq t_0$,
\begin{equation}\label{eq:ckn-relative-energy}
 \frac12\|w(t)\|_2^2+\nu\int_0^t\|\nabla w(s)\|_2^2\,ds
 \leq-\int_0^t\int_{\mathbb R^3}
 \sum_{i,j}w_iw_j\partial_j u_i\,dx\,ds
 \leq\int_0^t\kappa(s)\|w(s)\|_2^2\,ds.
\end{equation}
In the first step the extra term with $v_i=u_i+w_i$ vanishes because
$\sum_{i,j}u_iw_j\partial_j u_i=w\cdot\nabla(|u|^2/2)$ and
$\nabla\cdot w=0$. The original forcing cancels exactly.
For an explicit zero-data Gronwall argument, put
$Y(t)=\|w(t)\|_2^2$ and
$Z(t)=2\int_0^t\kappa(s)Y(s)\,ds$.
Then $0\leq Y\leq Z$, $Z(0)=0$, and
$Z'\leq2\kappa Z$ almost everywhere.
The derivative of $Z(t)\exp(-2\int_0^t\kappa)$ is nonpositive,
so this nonnegative function is identically zero. Thus $v=u$ almost
everywhere on $(0,t_0)$. Taking a countable increasing sequence
$t_0\uparrow T$ proves
\begin{equation}\label{eq:ckn-agreement}
 v=u\quad\hbox{almost everywhere on }\mathbb R^3\times(0,T).
\end{equation}
No terminal strong trace at $T$ and no uniqueness beyond $T$ entered
this argument.

\begin{theorem}[First-time off-axis exclusion for the original forced equation]
\label{thm:ckn-terminal}
For the smooth decaying data, force defined through $S>T$, and classical
strong branch just specified, suppose that $u$ is axisymmetric on
$[0,T)$. Then every $x_0$ with
$r_0=(x_{01}^2+x_{02}^2)^{1/2}>0$ has $\rho>0$ such that
\begin{equation}\label{eq:ckn-backward-regular}
 \mathop{\rm ess\,sup}_{B_\rho(x_0)\times(T-\rho^2,T)}|u|<\infty,
 \qquad T-\rho^2>0.
\end{equation}
In particular a first singularity in this backward-local velocity sense
cannot lie at positive radius. The continuation used in the proof need
not be axisymmetric at or after $T$.
\end{theorem}
\begin{proof}
Use $(v,P)$ on $(0,S)$ and \eqref{eq:ckn-agreement}. If no $\rho$
satisfies \eqref{eq:ckn-backward-regular}, then $u$ is essentially
unbounded in every sufficiently small backward cylinder at $(x_0,T)$.
Preterminal axisymmetry and rotation invariance of cylinders imply the
same statement at every $(R_\theta x_0,T)$. If any of those points were
regular for $v$, an open neighborhood where $v$ is bounded would contain
a sufficiently small backward cylinder, contradicting
\eqref{eq:ckn-agreement}. Thus every point of the positive-radius
circle is in $S(v)$. By \eqref{eq:ckn-circle-positive} this has positive
parabolic one-dimensional measure, whereas Theorem~\ref{thm:ckn-external}
gives zero. This contradiction proves \eqref{eq:ckn-backward-regular}.
\end{proof}

For smooth force, the spatial interior regularity included in the
external input turns a bounded neighborhood for $v$ into bounds for
its spatial derivatives on smaller neighborhoods. Thus the same argument
also excludes a positive-radius orbit of preterminal spatial-derivative
singularities: an unbounded derivative along a sequence converging to an
orbit point would contradict those smaller-neighborhood bounds and
\eqref{eq:ckn-agreement}. This specifies the extra regularity fact used
when the incoming report describes vorticity blow-up, rather than
velocity blow-up. It is not legitimate to apply a velocity singular-set
theorem directly to an arbitrary Euler vorticity singularity.

The conclusion is local in space. For a compact set $K$ disjoint from
the axis, the neighborhoods in \eqref{eq:ckn-backward-regular} have a
finite subcover of $K$. The minimum of their positive time lengths and
the maximum of their finite bounds give one velocity bound near $T$
on $K$. An assertion about a supremum on an unbounded region additionally
requires control at spatial infinity. A compact toroidal concentration
region bounded away from the axis has no such escape-to-infinity issue.

\subsection{Why preterminal smoothness alone is an insufficient replacement}

The following exact examples describe the missing force hypothesis;
they are not candidates in the prescribed smooth-through-$T$ force class.
Let $s=x_1^2+x_2^2+x_3^2$, let
$W(x)=(-x_2,x_1,0)$, and choose
$\eta\in C_c^\infty((0,\infty))$ equal to one on an interval around
$r_0^2$, where $r_0>0$. Put $V(x)=\eta(s)W(x)$.
Direct differentiation gives
\begin{align}\label{eq:ckn-example-field}
 \nabla\cdot V&=-2x_1x_2\eta'(s)+2x_1x_2\eta'(s)=0,\nonumber\\
 (V\cdot\nabla)V&=-\eta(s)^2(x_1,x_2,0),\nonumber\\
 \Delta V&=(4s\eta''(s)+10\eta'(s))W(x).
\end{align}
For the second identity, $W\cdot\nabla s=0$ and
$(W\cdot\nabla)W=-(x_1,x_2,0)$.
For the third, $\Delta\eta(s)=4s\eta''(s)+6\eta'(s)$ and
$2\nabla\eta(s)\cdot\nabla W_i=4\eta'(s)W_i$.
The field is smooth, compactly supported, and axisymmetric in exactly
the sense \eqref{eq:ckn-axisymmetry}.
For any smooth scalar $a:[0,T)\to\mathbb R$, define
\begin{equation}\label{eq:ckn-example-solution}
 u(x,t)=a(t)V(x),\quad p(x,t)=0,\quad
 f(x,t)=a'(t)V(x)+a(t)^2(V\cdot\nabla)V(x)-\nu a(t)\Delta V(x).
\end{equation}
Substitution of $u_t=a'V$, $(u\cdot\nabla)u=a^2(V\cdot\nabla)V$,
and $\Delta u=a\Delta V$ proves the original forced NS equation exactly.
Every time slice has finite energy $a(t)^2\|V\|_2^2$.

Take first $a(t)=(T-t)^{-1}$. On the plateau of $\eta$,
\begin{equation}\label{eq:ckn-example-force}
 f(x,t)=(T-t)^{-2}\big[(-x_2,x_1,0)-(x_1,x_2,0)\big],
 \qquad |f(x,t)|=\sqrt2\,(x_1^2+x_2^2)^{1/2}(T-t)^{-2}.
\end{equation}
In a small spatial neighborhood of $(r_0,0,0)$ the radius is bounded
below by $r_0/2$, so $\int|f|^q$ up to $T$ diverges for every
$q\geq1/2$. The velocity blows up at the entire circle, with finite
energy at each $t<T$ but without a uniform energy bound. This exact
example proves that pointwise preterminal smoothness and finite
time-slice energies do not imply the force or uniform-energy hypotheses.

Take instead $a(t)=\sin((T-t)^{-1})$. Then
$\sup_{t<T}\|u(t)\|_2\leq\|V\|_2$ and
$\int_0^T\|\nabla u\|_2^2\leq T\|\nabla V\|_2^2$.
Nevertheless $u$ has no limit at a point of the plateau where $V\ne0$,
and the tangential part of its force there has magnitude
$r(T-t)^{-2}|\cos((T-t)^{-1})|$. Its time integral diverges, since the
substitution $z=(T-t)^{-1}$ gives
$\int^\infty|\cos z|\,dz=\infty$.
Thus even uniform energy and integrated dissipation do not manufacture
an admissible continuation of a force that is only smooth before $T$.
The terminal theorem above avoids both defects by using the actual
prescribed force through $S>T$ and a suitable solution for that force.

\paragraph{Precise integration conclusion.}
The checked addition is the circle exclusion, its explicit force and
viscosity maps, and the preterminal weak--strong argument. The incoming
sentence about obtaining an axisymmetric weak continuation can be
replaced by the stronger proved statement that any suitable energy
continuation suffices. The result excludes retention of a fixed
positive-radius axisymmetric singular orbit for the original classical
viscous equation in the stated class. It does not assert a disproof of
another geometry, nonaxisymmetric mechanisms, concentration approaching
the axis, or an altered exact bridge between the Euler and viscous
equations. No existence theorem for any such alternative is inferred.


